\documentclass[10pt,a4paper]{amsart}
\usepackage[margin=19mm]{geometry}
\usepackage{amsmath,amssymb,mathtools}
\usepackage[scr=rsfs,cal=euler]{mathalfa}
\usepackage{xfrac}
\usepackage[unicode,hidelinks,bookmarks=false]{hyperref}
\newcommand{\ie}{i.e.\ }
\newcommand{\cf}{cf.\ }
\newcommand{\resp}{respectively\ }
\newcommand{\Subsection}{\subsection}
\providecommand{\nobreakdash}{\nobreak\hbox{-}}
\setlength{\emergencystretch}{2em}
\multlinegap0pt
\usepackage{bm}
\usepackage{amssymb}
\usepackage{mathtools}
\usepackage{float}
\usepackage[normalem]{ulem}
\usepackage[shortlabels]{enumitem}
\usepackage{csquotes}
\newtheorem{theorem}{Theorem}
\usepackage{cleveref}
\crefname{theorem}{Theorem}{Theorems}
\newcommand{\MA}{\bm{A}}
\newcommand{\MM}{\bm{M}}
\newcommand{\MS}{\bm{S}}
\newcommand{\Vo}{\bm{o}}
\newcommand{\Vv}{\bm{v}}
\newcommand{\Vx}{\bm{x}}
\newcommand{\Vy}{\bm{y}}
\title{Accelerating MPGP-type Methods Through Preconditioning}
\author{Jakub Kružík \and David Horák$^{* \dag}$}
\date{Source: arXiv:2507.00617v2 (CC BY 4.0)}
\begin{document}
\maketitle

\noindent\emph{Source and licence.} Accelerating MPGP-type Methods Through Preconditioning. Version arXiv:2507.00617v2 (CC BY 4.0), distributed by arXiv under the Creative Commons Attribution 4.0 International licence, http://creativecommons.org/licenses/by/4.0/, as recorded in the arXiv licence metadata for this exact version and shown on the abstract page https://arxiv.org/abs/2507.00617v2. Full licence notice: \texttt{licenses/arxiv-2507.00617-CC-BY-4.0.txt}.\par\smallskip

\noindent\textit{Editorial scope.} Counted unit: the article's theorem theo:condsimple (source0.tex lines 906--912). The article's complete original proof of it is delivered with it (source0.tex lines 913--944, one \texttt{proof} environment of 32 lines). No mathematics is written, repaired or re-derived: the statement and the proof are the article's own lines, in the article's own notation, and no retained line is substituted or re-flowed.  Retained uncounted as the source-local context the proof needs: lines 720--732; lines 890--900.  Nothing was omitted from any retained span, so the package carries no omission record.  No retained line contains a citation, so the wrapper carries no bibliography and no bibliography entry.  No retained line contains a figure environment, an image-inclusion command or drawing code, so no image is delivered and the package declares no asset dependency.  Editorial formatting only, in addition: an AMS \texttt{amsart} preamble that restates the article's own declarations and macros used by the retained lines, namely the environments \texttt{theorem} on separate counters, together with the standard packages those lines need.  The retained environments print in this document as Theorem 1 in order of appearance, whereas the article prints the same items with the numbering of its own full document.  The complete native source of the article as distributed in the arXiv e-print for this exact version is bundled unchanged as \texttt{sources/180821/source0.tex}.
\begin{theorem}[The error of the approximate preconditioning in face]
  \label{theo:eigvals}
  The (TODO additive?) error of the approximate preconditioning in face compared to the standard preconditioning in face is given by
\begin{equation*}
  \MM_{\mathcal{F}\mathcal{F}}^{-1} \MM_{\mathcal{F}\mathcal{A}} ( \MM_{\mathcal{A}\mathcal{A}} - \MM_{\mathcal{A}\mathcal{F}} \MM_{\mathcal{F}\mathcal{F}}^{-1} \MM_{\mathcal{F}\mathcal{A}})^{-1} \MM_{\mathcal{A}\mathcal{F}}.
\end{equation*}

Moreover, the non-zero eigenvalues of the preconditioned operator $\MM^{-1}\MA$ correspond to the eigenvalues of $\MS^{-1}\MA_{\mathcal{FF}}$, which are
\begin{equation*}
  1 = \lambda_1 = \dots = \lambda_{m-r} \le \dots \le \lambda_m,
\end{equation*}
where $r=\text{rank}(\MM_{\mathcal{AF}})$ and $m$ is the size of the free set.
\end{theorem}

\begin{theorem}[Kantorovich inequality]
  \label{theo:kant}
  For any symmetric positive definite matrix $\MA \in \mathbb{R}^{n \times n}$, the following inequalities are valid:
\begin{equation*}
  1 \leq \frac{\left(\Vx^{T} \MA \Vx \right) \left(\Vx^{T} \MA^{-1} \Vx\right)}{\left(\Vx^{T} \Vx\right)^{2}} \leq \frac{\left( \kappa\left( \MA \right) +1  \right)^{2}}{4\kappa(\MA)} \quad \forall \Vx \in \mathbb{R}^{n}\setminus \left\{\Vo \right\}.
\end{equation*}
Moreover, the upper bound is sharp, i.e.,
\begin{equation*}
  \sup_{\Vx \neq \Vo} \frac{\left(\Vx^{T} \MA \Vx \right) \left(\Vx^{T} \MA^{-1} \Vx\right)}{\left(\Vx^{T} \Vx\right)^{2}} = \frac{\left( \kappa\left( \MA \right) +1  \right)^{2}}{4\kappa(\MA)}.
\end{equation*}
\end{theorem}

\begin{theorem}[Conditioning of the preconditioned operator]
  \label{theo:condsimple}
  Let $\overline{\MM} = \MA$ and the inner preconditioner be the full inverse of $\overline{\MM}$. Then
    \begin{equation*}
      \kappa_{eff}\left(\MM^{-1}\MA\right) \leq \frac{\left( \kappa\left( \MA \right) +1  \right)^{2}}{4\kappa(\MA)}.
    \end{equation*}
\end{theorem}

\begin{proof}
  The maximal eigenvalue $\lambda_{max}$ of the preconditioned operator $\MM^{-1}\MA$ is by \cref{theo:eigvals} equal to the maximal eigenvalue of $\MS^{-1}\MA_{\mathcal{FF}}$, i.e.,
  \begin{equation*}
    \lambda_{max} = \sup_{\Vv_{\mathcal{F}} \neq \Vo} \frac{\Vv_{\mathcal{F}}^{T}\MA_{\mathcal{FF}}\Vv_{\mathcal{F}}}{\Vv_{\mathcal{F}}^{T}\MS_{\mathcal{FF}}\Vv_{\mathcal{F}}}.
  \end{equation*}
  From the Cauchy--Schwarz inequality
  \begin{equation}
    \label{eq:cs}
    \left(\Vv_{\mathcal{F}}^{T}\Vv_{\mathcal{F}}\right)^{2} = \left(\Vv_{\mathcal{F}}^{T}\MS^{\frac{1}{2}}\MS^{-\frac{1}{2}}\Vv_{\mathcal{F}}\right)^{2} \leq \left(\Vv_{\mathcal{F}}^{T}\MS\Vv_{\mathcal{F}}\right) \left(\Vv_{\mathcal{F}}^{T}\MS^{-1}\Vv_{\mathcal{F}}\right),
  \end{equation}
  we have
  \begin{equation*}
    \frac{1}{\Vv_{\mathcal{F}}^{T}\MS\Vv_{\mathcal{F}}} \leq \frac{\Vv_{\mathcal{F}}^{T}\MS^{-1}\Vv_{\mathcal{F}}}{\left(\Vv_{\mathcal{F}}^{T}\Vv_{\mathcal{F}}\right)^{2}}.
  \end{equation*}
  Multiplying both sides by $\Vv_{\mathcal{F}}^{T}\MA_{\mathcal{FF}}\Vv_{\mathcal{F}}$, we have for all $\Vv_{\mathcal{F}} \neq \Vo$
  \begin{equation*}
    \frac{\Vv_{\mathcal{F}}^{T}\MA_{\mathcal{FF}}\Vv_{\mathcal{F}}}{\Vv_{\mathcal{F}}^{T}\MS_{\mathcal{FF}}\Vv_{\mathcal{F}}} \leq
    \frac{\left(\Vv_{\mathcal{F}}^{T}\MA_{\mathcal{FF}}\Vv_{\mathcal{F}}\right) \left(\Vv_{\mathcal{F}}^{T}\MS^{-1}\Vv_{\mathcal{F}}\right)}{\left(\Vv_{\mathcal{F}}^{T}\Vv_{\mathcal{F}}\right)^{2}}.
  \end{equation*}
  Using $\left(\MA^{-1}\right)_\mathcal{FF} = \MS^{-1}$ from the block inversion formula and zero-padding $\Vv_{\mathcal{F}}$ as $\Vy = \begin{pmatrix} \Vv_{\mathcal{F}}\\ \Vo \end{pmatrix} \in \mathbb{R}^{n} \setminus\left\{\Vo\right\}$, we have
  \begin{equation*}
    \frac{\left(\Vv_{\mathcal{F}}^{T}\MA_{\mathcal{FF}}\Vv_{\mathcal{F}}\right) \left(\Vv_{\mathcal{F}}^{T}\MS^{-1}\Vv_{\mathcal{F}}\right)}{\left(\Vv_{\mathcal{F}}^{T}\Vv_{\mathcal{F}}\right)^{2}} =
    \frac{\left(\Vy^{T}\MA\Vy\right) \left(\Vy^{T}\MA^{-1}\Vy\right)}{\left(\Vy^{T}\Vy\right)^{2}}.
  \end{equation*}
  Since $\Vy$ is from a subset of $\mathbb{R}^{n}\setminus\left\{\Vo\right\}$, we bound the last term using the Kantorovich inequality (\cref{theo:kant}) by omitting the restriction on the zero components of $\Vy$ so that for all $\Vx \in \mathbb{R}^{n} \setminus\left\{\Vo\right\}$:
  \begin{equation*}
    \frac{\left(\Vy^{T}\MA\Vy\right) \left(\Vy^{T}\MA^{-1}\Vy\right)}{\left(\Vy^{T}\Vy\right)^{2}}
    \leq \frac{\left(\Vx^{T}\MA\Vx\right) \left(\Vx^{T}\MA^{-1}\Vx\right)}{\left(\Vx^{T}\Vx\right)^{2}}
    \leq \frac{\left( \kappa\left( \MA \right) +1 \right)^{2}}{4\kappa(\MA)}.
  \end{equation*}
  By \cref{theo:eigvals}, the minimal non-zero eigenvalue of the preconditioned operator $\MM^{-1}\MA$ is greater than or equal to $1$, which completes the proof.
\end{proof}

\end{document}
